\input format.tex

\usepackage{graphicx}
\graphicspath{{cores/}}

\usepackage{environ}
\usepackage{colortbl,array,booktabs}
\usepackage{tabularx}

\colorlet{TablaBordeSuperior}{topcolor}
\colorlet{TablaBordeInferior}{topcolor}
\colorlet{TablaCentroSuperior}{blue!1}
\colorlet{TablaCentroInferior}{blue!20}
\colorlet{FuenteCabeceraTabla}{white}

\newcolumntype{M}[1]{>{\centering\arraybackslash}m{#1}}
%\newcommand{\tabularxcolumn}[1]{>\arraybackslash}m{#1}}

\tcbset{rtab/.style={
freelance,
frame code={
 \path[top color=topcolor,bottom color=topcolor]
   ([yshift=-#1*(\baselineskip+2pt)]interior.north west) --
   ([yshift=-#1*(\baselineskip+2pt)]interior.north east) {[sharp corners]--
    ([yshift=3pt]interior.north east) --
    ([yshift=3pt]interior.north west)} -- cycle;

  },
interior code={},
 }
}

\newcommand\fuentecabecera[1]{\textcolor{black}{\textbf{#1}}}

\begin{document}

\vspace*{5mm}
%% 各章节
\setlength{\arrayrulewidth}{.2pt}
\fontsize{8.8pt}{11pt}\selectfont
\color{gray2}

\begin{center}
\fontsize{16pt}{\baselineskip}\selectfont\color{black70}\bf {三、详细检测结果}
\end{center}

\noindent{\fontsize{11pt}{11pt}\selectfont\color{black70}\bf {胖菌检测结果}}

\begin{tctabularxnew}{tabularx={m{4.8cm}<{\centering}m{5.2cm}<{\centering}m{4.52cm}<{\centering}}}
&&
\\[-6pt]
\cellcolor{topcolor} \raisebox{2.614pt}{\color{white}\bf 检测项目} &
\cellcolor{topcolor} \raisebox{2.614pt}{\color{white}\bf 健康状态提示} &
\cellcolor{topcolor} \raisebox{2.614pt}{\color{white}\bf 检测结果评价}
\end{tctabularxnew}

\vspace*{-4.25mm}
\fontsize{8pt}{11pt}\selectfont
\arrayrulecolor{gray2}
\begin{longtable}{m{4.8cm}m{5.2cm}<{\centering}m{0cm}@{}m{4.61cm}<{\centering}}
\hline
\parbox[c]{\hsize}{\vskip7pt {\lantxh 肠杆菌属\\如阴沟肠杆菌等，产生内毒素引起炎症，导致肥胖} \vskip7pt} & \parbox[c]{\hsize}{\vskip7pt\centerline{\raisebox{-1.5ex}{\includegraphics[width=5cm,height=0.55cm]{bar.pdf}}}\vskip7pt}  &
\hspace*{-4.83cm}\raisebox{-0.45ex}{\includegraphics{smile.pdf}}
 & \begin{minipage}{4.60cm}\begin{center}{
 {\lantxh 偏低}
  }\end{center} \end{minipage} \\
\hline
\parbox[c]{\hsize}{\vskip7pt {\lantxh 梭菌属\\多数为条件致病菌，如双酶梭菌等或其代谢物 可促进小肠中脂质的吸收} \vskip7pt} & \parbox[c]{\hsize}{\vskip7pt\centerline{\raisebox{-1.5ex}{\includegraphics[width=5cm,height=0.55cm]{bar.pdf}}}\vskip7pt}  &
\hspace*{-3.17cm}\raisebox{-0.45ex}{\includegraphics{cry.pdf}}
 & \begin{minipage}{4.60cm}\begin{center}{
 {\color{orange}\lantxh 偏高}
  }\end{center} \end{minipage} \\
\hline
\parbox[c]{\hsize}{\vskip7pt {\lantxh 布劳特氏菌属\\肥胖人群中布劳特氏菌属含量明显升高} \vskip7pt} & \parbox[c]{\hsize}{\vskip7pt\centerline{\raisebox{-1.5ex}{\includegraphics[width=5cm,height=0.55cm]{bar.pdf}}}\vskip7pt}  &
\hspace*{-3.17cm}\raisebox{-0.45ex}{\includegraphics{cry.pdf}}
 & \begin{minipage}{4.60cm}\begin{center}{
 {\color{orange}\lantxh 偏高}
  }\end{center} \end{minipage} \\
\hline
\parbox[c]{\hsize}{\vskip7pt {\lantxh 多尔氏菌属\\肠道里主要产气菌之一，在肥胖人群中易富集} \vskip7pt} & \parbox[c]{\hsize}{\vskip7pt\centerline{\raisebox{-1.5ex}{\includegraphics[width=5cm,height=0.55cm]{bar.pdf}}}\vskip7pt}  &
\hspace*{-3.17cm}\raisebox{-0.45ex}{\includegraphics{cry.pdf}}
 & \begin{minipage}{4.60cm}\begin{center}{
 {\color{orange}\lantxh 偏高}
  }\end{center} \end{minipage} \\
\hline
\parbox[c]{\hsize}{\vskip7pt {\lantxh 沙门氏菌属\\多数为致病菌，可能增加血液中脂多糖水平，促 进内脏脂肪、皮下脂肪炎症因子表达} \vskip7pt} & \parbox[c]{\hsize}{\vskip7pt\centerline{\raisebox{-1.5ex}{\includegraphics[width=5cm,height=0.55cm]{bar.pdf}}}\vskip7pt}  &
\hspace*{-4.83cm}\raisebox{-0.45ex}{\includegraphics{smile.pdf}}
 & \begin{minipage}{4.60cm}\begin{center}{
 {\lantxh 偏低}
  }\end{center} \end{minipage} \\
\hline
\parbox[c]{\hsize}{\vskip7pt {\lantxh 弯曲杆菌属\\多数为致病菌，肥胖患者中此类菌属含量增加} \vskip7pt} & \parbox[c]{\hsize}{\vskip7pt\centerline{\raisebox{-1.5ex}{\includegraphics[width=5cm,height=0.55cm]{bar.pdf}}}\vskip7pt}  &
\hspace*{-4.83cm}\raisebox{-0.45ex}{\includegraphics{smile.pdf}}
 & \begin{minipage}{4.60cm}\begin{center}{
 {\lantxh 偏低}
  }\end{center} \end{minipage} \\
\hline
\parbox[c]{\hsize}{\vskip7pt {\lantxh 志贺氏菌属\\是一类革兰阴性杆菌， 促进机体炎症因子的表达， 增加宿主肥胖风险} \vskip7pt} & \parbox[c]{\hsize}{\vskip7pt\centerline{\raisebox{-1.5ex}{\includegraphics[width=5cm,height=0.55cm]{bar.pdf}}}\vskip7pt}  &
\hspace*{-4.83cm}\raisebox{-0.45ex}{\includegraphics{smile.pdf}}
 & \begin{minipage}{4.60cm}\begin{center}{
 {\lantxh 偏低}
  }\end{center} \end{minipage} \\
\hline
\parbox[c]{\hsize}{\vskip7pt {\lantxh 克雷伯氏菌属\\多为致病菌，产生肠毒素引起慢性炎症，易引起肥胖} \vskip7pt} & \parbox[c]{\hsize}{\vskip7pt\centerline{\raisebox{-1.5ex}{\includegraphics[width=5cm,height=0.55cm]{bar.pdf}}}\vskip7pt}  &
\hspace*{-4.83cm}\raisebox{-0.45ex}{\includegraphics{smile.pdf}}
 & \begin{minipage}{4.60cm}\begin{center}{
 {\lantxh 偏低}
  }\end{center} \end{minipage} \\
\hline
\parbox[c]{\hsize}{\vskip7pt {\lantxh 脱硫弧菌属\\产生硫化氢，刺激肠道产生炎症反应，不利于肠道健康} \vskip7pt} & \parbox[c]{\hsize}{\vskip7pt\centerline{\raisebox{-1.5ex}{\includegraphics[width=5cm,height=0.55cm]{bar.pdf}}}\vskip7pt}  &
\hspace*{-3.17cm}\raisebox{-0.45ex}{\includegraphics{cry.pdf}}
 & \begin{minipage}{4.60cm}\begin{center}{
 {\color{orange}\lantxh 偏高}
  }\end{center} \end{minipage} \\
\hline
\parbox[c]{\hsize}{\vskip7pt {\lantxh 嗜胆菌属\\肥胖人群中此菌含量高，可能与长期高脂高蛋白饮食有关} \vskip7pt} & \parbox[c]{\hsize}{\vskip7pt\centerline{\raisebox{-1.5ex}{\includegraphics[width=5cm,height=0.55cm]{bar.pdf}}}\vskip7pt}  &
\hspace*{-4.83cm}\raisebox{-0.45ex}{\includegraphics{smile.pdf}}
 & \begin{minipage}{4.60cm}\begin{center}{
 {\lantxh 偏低}
  }\end{center} \end{minipage} \\
\hline
\end{longtable}

\clearpage

\vspace*{8mm}
%% 各章节
\setlength{\arrayrulewidth}{.2pt}
\fontsize{8.8pt}{11pt}\selectfont
\color{gray2}

\noindent{\fontsize{11pt}{11pt}\selectfont\color{black70}\bf {瘦菌检测结果}}

\begin{tctabularxnew}{tabularx={m{4.8cm}<{\centering}m{5.2cm}<{\centering}m{4.52cm}<{\centering}}}
&&
\\[-6pt]
\cellcolor{topcolor} \raisebox{2.614pt}{\color{white}\bf 检测项目} &
\cellcolor{topcolor} \raisebox{2.614pt}{\color{white}\bf 健康状态提示} &
\cellcolor{topcolor} \raisebox{2.614pt}{\color{white}\bf 检测结果评价}
\end{tctabularxnew}

\vspace*{-4.25mm}
\fontsize{8pt}{11pt}\selectfont
\arrayrulecolor{gray2}
\begin{longtable}{m{4.8cm}m{5.2cm}<{\centering}m{0cm}@{}m{4.61cm}<{\centering}}
\hline
\parbox[c]{\hsize}{\vskip7pt {\lantxh 阿克曼氏菌属\\降解粘蛋白、调节免疫，有利于肠黏膜完整性，保持正常体重} \vskip7pt} & \parbox[c]{\hsize}{\vskip7pt\centerline{\raisebox{-1.5ex}{\includegraphics[width=5cm,height=0.55cm]{bar.pdf}}}\vskip7pt}  &
\hspace*{-3.17cm}\raisebox{-0.45ex}{\includegraphics{cry.pdf}}
 & \begin{minipage}{4.60cm}\begin{center}{
 {\color{orange}\lantxh 偏低}
  }\end{center} \end{minipage} \\
\hline
\parbox[c]{\hsize}{\vskip7pt {\lantxh 双歧杆菌属\\有益菌，可以调节脂肪组织基因表达，可抑制系统性肥胖} \vskip7pt} & \parbox[c]{\hsize}{\vskip7pt\centerline{\raisebox{-1.5ex}{\includegraphics[width=5cm,height=0.55cm]{bar.pdf}}}\vskip7pt}  &
\hspace*{-3.17cm}\raisebox{-0.45ex}{\includegraphics{cry.pdf}}
 & \begin{minipage}{4.60cm}\begin{center}{
 {\color{orange}\lantxh 偏低}
  }\end{center} \end{minipage} \\
\hline
\parbox[c]{\hsize}{\vskip7pt {\lantxh 乳酸杆菌属\\有益菌，有利于降低胆固醇，减少产生内毒素，降低宿主肥胖风险} \vskip7pt} & \parbox[c]{\hsize}{\vskip7pt\centerline{\raisebox{-1.5ex}{\includegraphics[width=5cm,height=0.55cm]{bar.pdf}}}\vskip7pt}  &
\hspace*{-4.83cm}\raisebox{-0.45ex}{\includegraphics{smile.pdf}}
 & \begin{minipage}{4.60cm}\begin{center}{
 {\lantxh 偏高}
  }\end{center} \end{minipage} \\
\hline
\parbox[c]{\hsize}{\vskip7pt {\lantxh 克里斯滕森菌属\\与人体BMI呈负相关} \vskip7pt} & \parbox[c]{\hsize}{\vskip7pt\centerline{\raisebox{-1.5ex}{\includegraphics[width=5cm,height=0.55cm]{bar.pdf}}}\vskip7pt}  &
\hspace*{-3.17cm}\raisebox{-0.45ex}{\includegraphics{cry.pdf}}
 & \begin{minipage}{4.60cm}\begin{center}{
 {\lantxh 偏高}
  }\end{center} \end{minipage} \\
\hline
\parbox[c]{\hsize}{\vskip7pt {\lantxh 柔嫩梭菌属\\研究发现肠道中该类细菌的减少，会增加肥胖、炎 症性肠病等患病风险} \vskip7pt} & \parbox[c]{\hsize}{\vskip7pt\centerline{\raisebox{-1.5ex}{\includegraphics[width=5cm,height=0.55cm]{bar.pdf}}}\vskip7pt}  &
\hspace*{-3.17cm}\raisebox{-0.45ex}{\includegraphics{cry.pdf}}
 & \begin{minipage}{4.60cm}\begin{center}{
 {\color{orange}\lantxh 偏低}
  }\end{center} \end{minipage} \\
\hline
\parbox[c]{\hsize}{\vskip7pt {\lantxh 颤螺菌属\\帮助抗性淀粉和脂肪消化，保持正常体重，抑制肠道炎症} \vskip7pt} & \parbox[c]{\hsize}{\vskip7pt\centerline{\raisebox{-1.5ex}{\includegraphics[width=5cm,height=0.55cm]{bar.pdf}}}\vskip7pt}  &
\hspace*{-4.83cm}\raisebox{-0.45ex}{\includegraphics{smile.pdf}}
 & \begin{minipage}{4.60cm}\begin{center}{
 {\lantxh 正常}
  }\end{center} \end{minipage} \\
\hline
\parbox[c]{\hsize}{\vskip7pt {\lantxh 毛螺菌属\\发酵多种糖类产生短链脂肪酸，保护肠黏膜，抑制炎症，减轻肥胖症状} \vskip7pt} & \parbox[c]{\hsize}{\vskip7pt\centerline{\raisebox{-1.5ex}{\includegraphics[width=5cm,height=0.55cm]{bar.pdf}}}\vskip7pt}  &
\hspace*{-4.83cm}\raisebox{-0.45ex}{\includegraphics{smile.pdf}}
 & \begin{minipage}{4.60cm}\begin{center}{
 {\lantxh 正常}
  }\end{center} \end{minipage} \\
\hline
\parbox[c]{\hsize}{\vskip7pt {\lantxh 毛杆菌属\\肠道共生菌，发酵葡萄糖产生乳酸及少量乙酸和丁酸} \vskip7pt} & \parbox[c]{\hsize}{\vskip7pt\centerline{\raisebox{-1.5ex}{\includegraphics[width=5cm,height=0.55cm]{bar.pdf}}}\vskip7pt}  &
\hspace*{-3.17cm}\raisebox{-0.45ex}{\includegraphics{cry.pdf}}
 & \begin{minipage}{4.60cm}\begin{center}{
 {\color{orange}\lantxh 偏低}
  }\end{center} \end{minipage} \\
\hline
\end{longtable}

\end{document}

